% figures/w02-antiwindup-two.svg — 두 안티와인드업이 적분기 입력에서 갈라진다.
%
%   bash _tools/tikz2svg.sh figures/src/w02-antiwindup-two.tex figures/w02-antiwindup-two.svg
%
% 두 방식은 루프의 나머지가 전부 같다. 다른 것은 **적분기 앞에서 무엇이 일어나는가**
% 하나뿐이므로, 그 자리만 나란히 놓고 나머지는 같은 모양으로 그린다.
% 전체 루프(포화 뒤가 왜 끊기는가)는 w02-windup.svg 가 이미 그렸다 — 여기서
% 선체와 바깥 되먹임을 다시 그리지 않는 이유다.
%
% 폭: 두 판 합쳐 약 172 mm. A4 본문 폭 186 mm 안에 들어간다 (CLAUDE.md 4 규칙 6-0).
\documentclass[border=5pt]{standalone}
% gnc-style.tex — 이 볼트의 TikZ 그림이 공유하는 스타일.
%
% 저널 그림 규약을 스타일 하나로 굳혀 둔다. 그림마다 선 굵기나 화살촉을
% 다시 정하지 않는다 — 그것이 그림이 제각각으로 보이는 원인이다.
%
% 쓰는 법:  % gnc-style.tex — 이 볼트의 TikZ 그림이 공유하는 스타일.
%
% 저널 그림 규약을 스타일 하나로 굳혀 둔다. 그림마다 선 굵기나 화살촉을
% 다시 정하지 않는다 — 그것이 그림이 제각각으로 보이는 원인이다.
%
% 쓰는 법:  % gnc-style.tex — 이 볼트의 TikZ 그림이 공유하는 스타일.
%
% 저널 그림 규약을 스타일 하나로 굳혀 둔다. 그림마다 선 굵기나 화살촉을
% 다시 정하지 않는다 — 그것이 그림이 제각각으로 보이는 원인이다.
%
% 쓰는 법:  \input{gnc-style}  (figures/src/ 안에서)

\usepackage{tikz}
\usepackage{amsmath}
\usepackage{xcolor}
\usetikzlibrary{arrows.meta, positioning, calc, fit, backgrounds, decorations.pathmorphing}

% ---- 색 --------------------------------------------------------------------
% 색은 강조 하나에만 쓴다. 나머지는 검정.
\definecolor{ink}{HTML}{1A1A1A}
\definecolor{accent}{HTML}{7C3AED}
\definecolor{warn}{HTML}{B91C1C}
\definecolor{soft}{HTML}{666666}

% ---- 선과 화살촉 -----------------------------------------------------------
\tikzset{
  % 신호선. 화살촉은 작고 뾰족한 것 하나뿐이다.
  sig/.style   = {ink, line width=0.5pt, -{Stealth[length=4pt, width=3pt]}},
  wire/.style  = {ink, line width=0.5pt},
  % 강조 경로 — 파선으로 구분한다. 흑백 인쇄에서도 살아야 하기 때문이다.
  asig/.style  = {accent, line width=0.5pt, densely dashed,
                  -{Stealth[length=4pt, width=3pt]}},
  awire/.style = {accent, line width=0.5pt, densely dashed},
  % 블록. 흰 바탕, 직각, 검은 테두리.
  blk/.style   = {draw=ink, line width=0.55pt, fill=white,
                  minimum width=17mm, minimum height=8mm, inner sep=2pt, font=\small},
  % 합산점
  sum/.style   = {draw=ink, line width=0.55pt, fill=white, circle,
                  inner sep=0pt, minimum size=4.6mm, font=\scriptsize},
  asum/.style  = {draw=accent, line width=0.55pt, fill=white, circle,
                  inner sep=0pt, minimum size=4.6mm, font=\scriptsize},
  % 분기점
  dot/.style   = {ink, fill=ink, circle, inner sep=0pt, minimum size=1.5mm},
  % 신호 이름 — 선 위에 작게
  sname/.style = {font=\small, inner sep=1.5pt},
  % 그림 안의 짧은 설명. 그림 안의 글은 최소로 둔다
  note/.style  = {font=\scriptsize, soft, inner sep=1.5pt},
  anote/.style = {font=\scriptsize, accent, inner sep=1.5pt},
  % 축
  axis/.style  = {ink, line width=0.5pt},
  tick/.style  = {font=\scriptsize, soft},
}

% 캡션 한 줄 — 그림 아래 가는 선과 함께
\newcommand{\gnccaption}[2]{%
  \node[anchor=north west, text width=#1, align=left, font=\scriptsize, ink]
        at (current bounding box.south west) {\vspace{2pt}#2};}
  (figures/src/ 안에서)

\usepackage{tikz}
\usepackage{amsmath}
\usepackage{xcolor}
\usetikzlibrary{arrows.meta, positioning, calc, fit, backgrounds, decorations.pathmorphing}

% ---- 색 --------------------------------------------------------------------
% 색은 강조 하나에만 쓴다. 나머지는 검정.
\definecolor{ink}{HTML}{1A1A1A}
\definecolor{accent}{HTML}{7C3AED}
\definecolor{warn}{HTML}{B91C1C}
\definecolor{soft}{HTML}{666666}

% ---- 선과 화살촉 -----------------------------------------------------------
\tikzset{
  % 신호선. 화살촉은 작고 뾰족한 것 하나뿐이다.
  sig/.style   = {ink, line width=0.5pt, -{Stealth[length=4pt, width=3pt]}},
  wire/.style  = {ink, line width=0.5pt},
  % 강조 경로 — 파선으로 구분한다. 흑백 인쇄에서도 살아야 하기 때문이다.
  asig/.style  = {accent, line width=0.5pt, densely dashed,
                  -{Stealth[length=4pt, width=3pt]}},
  awire/.style = {accent, line width=0.5pt, densely dashed},
  % 블록. 흰 바탕, 직각, 검은 테두리.
  blk/.style   = {draw=ink, line width=0.55pt, fill=white,
                  minimum width=17mm, minimum height=8mm, inner sep=2pt, font=\small},
  % 합산점
  sum/.style   = {draw=ink, line width=0.55pt, fill=white, circle,
                  inner sep=0pt, minimum size=4.6mm, font=\scriptsize},
  asum/.style  = {draw=accent, line width=0.55pt, fill=white, circle,
                  inner sep=0pt, minimum size=4.6mm, font=\scriptsize},
  % 분기점
  dot/.style   = {ink, fill=ink, circle, inner sep=0pt, minimum size=1.5mm},
  % 신호 이름 — 선 위에 작게
  sname/.style = {font=\small, inner sep=1.5pt},
  % 그림 안의 짧은 설명. 그림 안의 글은 최소로 둔다
  note/.style  = {font=\scriptsize, soft, inner sep=1.5pt},
  anote/.style = {font=\scriptsize, accent, inner sep=1.5pt},
  % 축
  axis/.style  = {ink, line width=0.5pt},
  tick/.style  = {font=\scriptsize, soft},
}

% 캡션 한 줄 — 그림 아래 가는 선과 함께
\newcommand{\gnccaption}[2]{%
  \node[anchor=north west, text width=#1, align=left, font=\scriptsize, ink]
        at (current bounding box.south west) {\vspace{2pt}#2};}
  (figures/src/ 안에서)

\usepackage{tikz}
\usepackage{amsmath}
\usepackage{xcolor}
\usetikzlibrary{arrows.meta, positioning, calc, fit, backgrounds, decorations.pathmorphing}

% ---- 색 --------------------------------------------------------------------
% 색은 강조 하나에만 쓴다. 나머지는 검정.
\definecolor{ink}{HTML}{1A1A1A}
\definecolor{accent}{HTML}{7C3AED}
\definecolor{warn}{HTML}{B91C1C}
\definecolor{soft}{HTML}{666666}

% ---- 선과 화살촉 -----------------------------------------------------------
\tikzset{
  % 신호선. 화살촉은 작고 뾰족한 것 하나뿐이다.
  sig/.style   = {ink, line width=0.5pt, -{Stealth[length=4pt, width=3pt]}},
  wire/.style  = {ink, line width=0.5pt},
  % 강조 경로 — 파선으로 구분한다. 흑백 인쇄에서도 살아야 하기 때문이다.
  asig/.style  = {accent, line width=0.5pt, densely dashed,
                  -{Stealth[length=4pt, width=3pt]}},
  awire/.style = {accent, line width=0.5pt, densely dashed},
  % 블록. 흰 바탕, 직각, 검은 테두리.
  blk/.style   = {draw=ink, line width=0.55pt, fill=white,
                  minimum width=17mm, minimum height=8mm, inner sep=2pt, font=\small},
  % 합산점
  sum/.style   = {draw=ink, line width=0.55pt, fill=white, circle,
                  inner sep=0pt, minimum size=4.6mm, font=\scriptsize},
  asum/.style  = {draw=accent, line width=0.55pt, fill=white, circle,
                  inner sep=0pt, minimum size=4.6mm, font=\scriptsize},
  % 분기점
  dot/.style   = {ink, fill=ink, circle, inner sep=0pt, minimum size=1.5mm},
  % 신호 이름 — 선 위에 작게
  sname/.style = {font=\small, inner sep=1.5pt},
  % 그림 안의 짧은 설명. 그림 안의 글은 최소로 둔다
  note/.style  = {font=\scriptsize, soft, inner sep=1.5pt},
  anote/.style = {font=\scriptsize, accent, inner sep=1.5pt},
  % 축
  axis/.style  = {ink, line width=0.5pt},
  tick/.style  = {font=\scriptsize, soft},
}

% 캡션 한 줄 — 그림 아래 가는 선과 함께
\newcommand{\gnccaption}[2]{%
  \node[anchor=north west, text width=#1, align=left, font=\scriptsize, ink]
        at (current bounding box.south west) {\vspace{2pt}#2};}


% 이 그림에서만 쓰는 좁은 블록 — 두 판을 나란히 놓아야 하므로 기본 17mm 은 넓다
\tikzset{
  nblk/.style = {draw=ink, line width=0.55pt, fill=white,
                 minimum width=11mm, minimum height=7.5mm, inner sep=1.5pt, font=\small},
  ablk/.style = {draw=accent, line width=0.7pt, fill=white,
                 minimum width=13mm, minimum height=8.5mm, inner sep=1.5pt,
                 font=\small, text=accent},
}

\begin{document}
\begin{tikzpicture}[x=1mm, y=1mm]

\node[font=\small\bfseries, anchor=north west] at (-12,34)
     {Two anti-windup schemes, and the one place they differ};
\node[anchor=north west, align=left, text width=172mm, font=\scriptsize, soft] at (-12,28) {%
  Everything outside the violet block is identical, including the excess
  $\varepsilon = X_{\text{cmd}} - X_{\text{sat}}$ that both schemes measure. The schemes
  disagree only about what to do with it at the integrator's input.};

% =======================================================================
% 두 판을 같은 매크로로 그린다. \PX 가 판의 원점.
% =======================================================================

% ---------- (a) clamping ----------------------------------------------
\begin{scope}[shift={(0,0)}]
  \node[font=\small\bfseries, anchor=west, text=ink] at (-10,14) {(a)\quad clamping};
  \node[note, anchor=west, text width=76mm] at (-10,10)
       {switch the integrator \emph{off} while it would wind further in};

  \node[nblk] (kp)  at (8,4)    {$K_p$};
  \node[nblk] (ki)  at (8,-8)   {$K_i$};
  \node[ablk] (gate) at (26,-8) {};
  \node[nblk] (int) at (44,-8)  {$1/s$};
  \node[sum]  (s2)  at (58,0)   {};
  \node[nblk] (sat) at (70,0)   {};

  % 스위치 글리프 — 열린 접점
  \draw[accent, line width=0.6pt] ([shift={(-4mm,0)}]gate.center) -- ([shift={(-1.6mm,0)}]gate.center);
  \draw[accent, line width=0.6pt] ([shift={( 1.6mm,0)}]gate.center) -- ([shift={( 4mm,0)}]gate.center);
  \draw[accent, line width=0.7pt] ([shift={(-1.6mm,0)}]gate.center) -- ([shift={( 1.3mm,2.3mm)}]gate.center);
  \fill[accent] ([shift={(-1.6mm,0)}]gate.center) circle (0.42);
  \fill[accent] ([shift={( 1.6mm,0)}]gate.center) circle (0.42);

  % 포화 글리프
  \draw[wire] ([shift={(-4.5mm,-2.1mm)}]sat.center) -- ([shift={(-1.7mm,-2.1mm)}]sat.center)
              -- ([shift={( 1.7mm, 2.1mm)}]sat.center) -- ([shift={( 4.5mm, 2.1mm)}]sat.center);

  % 배선
  \coordinate (br) at (-2,-2);
  \draw[sig] (-10,-2) -- node[sname, above, pos=0.15] {$e$} (br);
  \node[dot] at (br) {};
  \draw[wire] (br) |- (kp.west);   \draw[wire] (br) |- (ki.west);
  \draw[sig]  (ki) -- (gate);
  \draw[sig]  (gate) -- (int);
  \draw[wire] (kp) -| (s2);
  \draw[wire] (int) -| (s2);
  %  X_cmd 는 선 위에 두지 않는다 — 합산점과 포화 블록 사이가 9mm 라 이름표가
  %  반드시 한쪽에 걸린다. 내려가는 탭 선 옆의 빈 곳에 붙인다.
  \draw[sig]  (s2) -- (sat);
  \draw[sig]  (sat) -- (81,0);
  \node[sname, anchor=west] at (81,0) {$X_{\text{sat}}$};

  % epsilon 을 재는 곳. 탭은 X_cmd 이름표와 겹치지 않게 오른쪽으로 물린다.
  \coordinate (ta) at (66,0);   \coordinate (tb) at (78,0);
  \node[dot] at (ta) {};  \node[dot] at (tb) {};
  \node[asum] (ep) at (72,-20) {};
  \node[anote, anchor=north] at (72,-23.5) {$\varepsilon$};
  \draw[awire] (ta) |- (ep);
  \draw[awire] (tb) |- (ep);
  \node[anote, anchor=east] at (69.4,-17.6) {$+$};
  \node[anote, anchor=west] at (74.6,-17.6) {$-$};
  \node[sname, anchor=east] at (65.2,-11) {$X_{\text{cmd}}$};

  % 게이트로 가는 논리
  \draw[asig] (ep) -| (26,-14.5);
  \node[anote, anchor=north, align=center, text width=54mm] at (26,-24)
       {open when $\varepsilon \neq 0$ \textbf{and} $\operatorname{sign}(e)=\operatorname{sign}(\varepsilon)$,\\
        so that $\dot I = 0$ — the memory is frozen};
\end{scope}

% ---------- (b) back-calculation --------------------------------------
\begin{scope}[shift={(104,0)}]
  \node[font=\small\bfseries, anchor=west, text=ink] at (-10,14) {(b)\quad back-calculation};
  \node[note, anchor=west, text width=76mm] at (-10,10)
       {feed the excess \emph{back} so the integrator is pulled to the limit};

  \node[nblk] (kp)  at (8,4)    {$K_p$};
  \node[nblk] (ki)  at (8,-8)   {$K_i$};
  \node[asum] (bs)  at (26,-8)  {};
  \node[nblk] (int) at (44,-8)  {$1/s$};
  \node[sum]  (s2)  at (58,0)   {};
  \node[nblk] (sat) at (70,0)   {};
  \node[ablk] (kaw) at (26,-20) {$K_{\text{aw}}$};

  \draw[wire] ([shift={(-4.5mm,-2.1mm)}]sat.center) -- ([shift={(-1.7mm,-2.1mm)}]sat.center)
              -- ([shift={( 1.7mm, 2.1mm)}]sat.center) -- ([shift={( 4.5mm, 2.1mm)}]sat.center);

  \coordinate (br) at (-2,-2);
  \draw[sig] (-10,-2) -- node[sname, above, pos=0.15] {$e$} (br);
  \node[dot] at (br) {};
  \draw[wire] (br) |- (kp.west);   \draw[wire] (br) |- (ki.west);
  \draw[sig]  (ki) -- (bs);
  \draw[sig]  (bs) -- (int);
  \draw[wire] (kp) -| (s2);
  \draw[wire] (int) -| (s2);
  %  X_cmd 는 선 위에 두지 않는다 — (a) 판과 같은 이유다.
  \draw[sig]  (s2) -- (sat);
  \draw[sig]  (sat) -- (81,0);
  \node[sname, anchor=west] at (81,0) {$X_{\text{sat}}$};

  \coordinate (ta) at (66,0);   \coordinate (tb) at (78,0);
  \node[dot] at (ta) {};  \node[dot] at (tb) {};
  \node[asum] (ep) at (72,-20) {};
  \node[anote, anchor=north] at (72,-23.5) {$\varepsilon$};
  \draw[awire] (ta) |- (ep);
  \draw[awire] (tb) |- (ep);
  \node[anote, anchor=east] at (69.4,-17.6) {$+$};
  \node[anote, anchor=west] at (74.6,-17.6) {$-$};
  \node[sname, anchor=east] at (65.2,-11) {$X_{\text{cmd}}$};

  % epsilon -> Kaw -> 적분기 입력의 뺄셈
  \draw[asig] (ep) -- (kaw);
  \draw[asig] (kaw) -- (bs);
  \node[anote, anchor=east] at (24.6,-11.5) {$-$};
  \node[anote, anchor=north, align=center, text width=54mm] at (26,-26)
       {$\dot I = K_i e - K_{\text{aw}}\varepsilon$, a continuous pull\\
        towards the value that makes $X_{\text{cmd}} = X_{\text{sat}}$};
\end{scope}

\end{tikzpicture}
\end{document}
